\section{从多任务 RL 到元学习}

\subsection{回顾多任务 RL}

上一讲讨论了多任务 RL：通过将状态增广为 $(s, z)$（其中 $z$ 是任务描述），训练条件策略 $\pi(a|s,z)$ 来处理多个任务。但多任务 RL 有一个关键局限：它假设在\textbf{训练时就知道所有任务}。

\begin{figure}[H]
\centering
\includegraphics[width=0.9\textwidth]{slides-images/slide-03.jpg}
\caption{多任务 RL 回顾：把任务描述并入状态，统一成 joint MDP}
\end{figure}

\begin{knowledgebox}{多任务 RL 的局限}
多任务 RL 通过一个共享策略覆盖所有任务。如果测试时遇到的新任务与训练任务差异较大，性能可能急剧下降。此外，随着任务数量增多，学习一个对所有任务都表现良好的策略变得越来越困难（任务干扰问题）。
\end{knowledgebox}

\begin{figure}[H]
\centering
\includegraphics[width=0.9\textwidth]{slides-images/slide-04.jpg}
\caption{Goal-conditioned learning 是多任务 RL 的一个特殊情形}
\end{figure}

\subsection{元学习的核心思想}

\begin{importantbox}{Meta-RL 的定义}
元强化学习（Meta-RL）的目标不是学习"如何解决一个具体任务"，而是\textbf{学习"如何快速学习新任务"}。

形式化地说：给定一个任务分布 $p(\mathcal{T})$，Meta-RL 在训练阶段利用从 $p(\mathcal{T})$ 采样的大量任务学习一个"学习算法"或"快速适应机制"。测试时面对一个新任务 $\mathcal{T}_{\text{new}} \sim p(\mathcal{T})$，仅需少量数据即可适应。
\end{importantbox}

\subsection{为什么需要 Meta-RL}

讲者用一个非常生活化的例子来说明问题：学会操作一台新的 espresso machine。若把这个任务交给从零开始训练的 RL agent，可能要数百万次尝试；而人类只要有一点相关经验和一些指令，就能在几分钟内上手。差异不在于人类更会 ``优化奖励''，而在于人类\textbf{从不从零开始}。

\begin{figure}[H]
\centering
\includegraphics[width=0.9\textwidth]{slides-images/slide-06.jpg}
\caption{元强化学习的动机：面对新任务时，人类会利用过往经验快速适应}
\end{figure}

\begin{knowledgebox}{Meta-RL 试图弥补的能力缺口}
普通 RL 假定每个新任务都要重新收集大量交互数据；Meta-RL 则试图把 ``过去做过相似任务'' 这一事实编码进学习系统，使其能在新任务上更快进入有效行为区间。
\end{knowledgebox}

\subsection{从 transfer learning 到 meta-learning}

这节课还把几种常见迁移范式摆在一起比较：
\begin{enumerate}
    \item \textbf{Forward transfer}：先做 source task，再 fine-tune 到 target task。
    \item \textbf{Multi-task transfer}：同时在多个任务上训练，用 task descriptor 支持零样本泛化。
    \item \textbf{Meta-learning}：在训练阶段就显式优化 ``如何快速适应新任务''。
\end{enumerate}

\begin{figure}[H]
\centering
\includegraphics[width=0.9\textwidth]{slides-images/slide-07.jpg}
\caption{Transfer、multi-task 与 meta-learning 的问题设定差异}
\end{figure}

\begin{warningbox}{Meta-RL 的泛化边界}
Meta-RL 并不承诺能适应任意新任务。它依赖一个关键假设：meta-test 时遇到的新任务，仍然来自与 meta-train 相近的任务分布。如果测试任务完全跳出训练分布，适应速度与最终表现都可能显著下降。
\end{warningbox}

\subsection{Few-shot 学习视角}

课程还借助 few-shot image classification 与 LLM 的 in-context learning 做了类比。核心共同点是：模型接收的不只是输入 $x$，还接收一小组示例 $D_{\text{train}}$，并据此推断当前问题属于哪种任务、应该如何输出。

\begin{figure}[H]
\centering
\includegraphics[width=0.9\textwidth]{slides-images/slide-09.jpg}
\caption{Few-shot learning 的抽象：预测函数同时依赖输入和少量训练示例}
\end{figure}

\subsection{本章小结}

Meta-RL 将学习的层次从"学习策略"提升到"学习如何学习策略"，适用于需要快速适应新任务的场景。

